\documentclass{article}

\usepackage{neurips_2019_author_response}

\usepackage[utf8]{inputenc} % allow utf-8 input
\usepackage[T1]{fontenc}    % use 8-bit T1 fonts
\usepackage[
  colorlinks=true,
]{hyperref}
\usepackage{url}            % simple URL typesetting
\usepackage{booktabs}       % professiona\textbf{L}-quality tables
\usepackage{amsfonts}       % blackboard math symbols
\usepackage{nicefrac}       % compact symbols for 1/2, etc.
\usepackage{microtype}      % microtypography
\usepackage{lipsum}

% extra package and macros
\usepackage{xcolor}
\usepackage{amssymb}% http://ctan.org/pkg/amssymb
\usepackage{pifont}% http://ctan.org/pkg/pifont
\usepackage{graphicx}
\newcommand{\cmark}{\ding{51}}%
\newcommand{\xmark}{\ding{55}}%
\usepackage{multirow}
\usepackage{wrapfig,lipsum,booktabs}

% %% layout
\newcolumntype{L}[1]{>{\raggedright\let\newline\\\arraybackslash\hspace{0pt}}m{#1}}
\newcolumntype{C}[1]{>{\centering\let\newline\\\arraybackslash\hspace{0pt}}m{#1}}
\newcolumntype{R}[1]{>{\raggedleft\let\newline\\\arraybackslash\hspace{0pt}}m{#1}}

\newcommand{\xpar}[1]{\noindent\textbf{#1}\ \ }
\newcommand{\vpar}[1]{\vspace{3mm}\noindent\textbf{#1}\ \ }

%% notations
\newcommand{\sect}[1]{Section~\ref{#1}}
\newcommand{\sects}[1]{Sections~\ref{#1}}
\newcommand{\eqn}[1]{Equation~\ref{#1}}
\newcommand{\eqns}[1]{Equations~\ref{#1}}
\newcommand{\fig}[1]{Figure~\ref{#1}}
\newcommand{\figs}[1]{Figures~\ref{#1}}
\newcommand{\tab}[1]{Table~\ref{#1}}
\newcommand{\tabs}[1]{Tables~\ref{#1}}

\newcommand{\ignorethis}[1]{}
\newcommand{\norm}[1]{\lVert#1\rVert}
\newcommand{\fcseven}{$\mbox{fc}_7$}

\renewcommand*{\thefootnote}{\fnsymbol{footnote}}


\DeclareMathOperator*{\argmin}{arg\,min}
\DeclareMathOperator*{\argmax}{arg\,max}

\def\naive{na\"{\i}ve\xspace}
\def\Naive{Na\"{\i}ve\xspace}

\makeatletter
\DeclareRobustCommand\onedot{\futurelet\@let@token\@onedot}
\def\@onedot{\ifx\@let@token.\else.\null\fi\xspace}

\def\iid{\emph{i.i.d}\onedot}
\def\eg{\emph{e.g}\onedot} \def\Eg{\emph{E.g}\onedot}
\def\ie{\emph{i.e}\onedot} \def\Ie{\emph{I.e}\onedot}
\def\cf{\emph{c.f}\onedot} \def\Cf{\emph{C.f}\onedot}
\def\etc{\emph{etc}\onedot} \def\vs{\emph{vs}\onedot}
\def\wrt{w.r.t\onedot} \def\dof{d.o.f\onedot}
\def\etal{\emph{et al}\onedot}
\makeatother

%% comments
\definecolor{citecolor}{RGB}{34,139,34}
\definecolor{mydarkblue}{rgb}{0,0.08,1}
\definecolor{mydarkgreen}{rgb}{0.02,0.6,0.02}
\definecolor{mydarkred}{rgb}{0.8,0.02,0.02}
\definecolor{mydarkorange}{rgb}{0.40,0.2,0.02}
\definecolor{mypurple}{RGB}{111,0,255}
\definecolor{myred}{rgb}{1.0,0.0,0.0}
\definecolor{mygold}{rgb}{0.75,0.6,0.12}
\definecolor{mydarkgray}{rgb}{0.66, 0.66, 0.66}

\newcommand{\ligeng}[1]{\textcolor{mydarkred}{[Ligeng: #1]}}
\newcommand{\zhijian}[1]{\textcolor{mydarkblue}{[Zhijian: #1]}}
\newcommand{\SH}[1]{\textcolor{mydarkred}{[Song: #1]}}

\newcommand{\myparagraph}[1]{\paragraph{#1}}
\newcommand{\minisection}[1]{\vspace{5pt}\noindent\textbf{#1.}}

% extra package and macros
\usepackage{xcolor}
\usepackage{amssymb}% http://ctan.org/pkg/amssymb
\usepackage{pifont}% http://ctan.org/pkg/pifont
\usepackage{graphicx}
\newcommand{\cmark}{\ding{51}}%
\newcommand{\xmark}{\ding{55}}%
\usepackage{multirow}
\usepackage{wrapfig,lipsum,booktabs}
\AtBeginDocument{% optionally set colors to your liking
  \hypersetup{
    urlcolor=red,
    citecolor=green,
    linkcolor=red,
  }%
}
\definecolor{citecolor}{RGB}{34,139,34}
\definecolor{mydarkblue}{rgb}{0,0.2,0.8}
\definecolor{mydarkgreen}{rgb}{0.02,0.6,0.02}
\definecolor{mydarkred}{rgb}{0.8,0.02,0.02}
\newcommand{\ligeng}[1]{\textcolor{mydarkred}{[Ligeng: #1]}}
\newcommand{\zhijian}[1]{\textcolor{mydarkblue}{[Zhijian: #1]}}
\newcommand{\SH}[1]{\textcolor{mydarkred}{[Song: #1]}}
    
\makeatletter
\DeclareRobustCommand\onedot{\futurelet\@let@token\@onedot}
\def\@onedot{\ifx\@let@token.\else.\null\fi\xspace}

\def\iid{\emph{i.i.d}\onedot}
\def\eg{\emph{e.g}\onedot} \def\Eg{\emph{E.g}\onedot}
\def\ie{\emph{i.e}\onedot} \def\Ie{\emph{I.e}\onedot}
\def\cf{\emph{c.f}\onedot} \def\Cf{\emph{C.f}\onedot}
\def\etc{\emph{etc}\onedot} \def\vs{\emph{vs}\onedot}
\def\wrt{w.r.t\onedot} \def\dof{d.o.f\onedot}
\def\etal{\emph{et al}\onedot}
\makeatother

\newcommand{\revquote}[2]{
    % \vspace{0.02\linewidth}
    \noindent{\textcolor{mydarkblue}{\textbf{#1:} #2}}
}
\newcommand{\revquoteB}[4]{
    % \vspace{0.02\linewidth}
    \noindent{\small{\textcolor{mydarkblue}{\textbf{#1:} #2; \textbf{#3:} #4}}}
}

\newcommand{\lgcmark}{\textcolor{green}{\cmark}}
\newcommand{\lgxmark}{\textcolor{red}{\xmark}}

\usepackage{xstring}
\usepackage{amstext}
\usepackage{subcaption}
\usepackage{etoolbox}% http://ctan.org/pkg/etoolbox
\makeatletter
\newcommand{\lgcite}[1]{%
  \begingroup%
  % count the elements
  \@tempcnta=\z@
  \renewcommand*\do[1]{\advance\@tempcnta\@ne}%
  \docsvlist{#1}%
  \chardef\@listsize\@tempcnta%
  \@tempcnta=\@ne[%
  \renewcommand*\do[1]{%
    \ifnum\@tempcnta<\@listsize
      \textcolor{citecolor}{##1},\penalty\z@ \advance\@tempcnta\@ne%
    \else
      \unpenalty\textcolor{citecolor}{##1}]%
    \fi}%
  \docsvlist{#1}%
  \endgroup%
}
\makeatother


\begin{document}

    
We thank all reviewers for their comments. All reviewers think it is an interesting paper. 
%
R1's review summarizes our contribution well: ``DLG is the first to shows a malicious player can recover private training data in collaborative learning scenario.''
%
Both R1 and R3 are positive overall in their comments (R1 ``easy to read and well structured'', ``raises an important privacy issue'', R3 ``easy to read and understand'', ``it is surprising that obtaining the training datasets is possible by only utilizing the gradients'' ). For all typos/grammar mistakes, we have revised our writing accordingly. 

\revquote{R2}{DLG may not work for accumulated gradients / Contrived settings.}
% 
This is a misunderstanding: DLG is still effective in federated learning (Tab.~\ref{tab:accumulation}).
% 
In the real-world case, a common workflow is to firstly deliver a pre-trained model to users' devices and fine-tune it by Federated Learning\footnote{Towards federated learning at scale:  System design~\url{https://arxiv.org/abs/1902.01046}}. In this case, the gradient and learning rate are both small, thus the weight changes are small too. Thereby it can be approximated as multi-batch case and this is still possible to attack.
% 
Nowadays, noisy / sparse / accumulated gradients are just 
\textit{optional} choices for training acceleration, but actually they are 
\textit{essential} techniques to protect the training set. \textbf{Our work aims to raise people's awareness about the security of gradients}.

\begin{minipage}{0.6\linewidth}
    \centering
    \small{
        \begin{tabular}{c|c|c|c|c}
                                  &   Iterations=1 &  Iterations=2 &  Iterations=3 &  Iterations=4 \\ \hline
             MSE   &  $3.3 \times 10^{-6}$  & $3.5 \times 10^{-3}$ & $3.0 \times 10^{-3}$ &  $1.8 \times 10^{-2}$ \\ 
        \end{tabular}
    \captionof{table}{The effectiveness of DLG on federated learning for different communication frequency.}
    \label{tab:accumulation}
    }
\end{minipage}
% 
\begin{minipage}{0.4\linewidth}
    \centering
    \small{
        \begin{tabular}{c|c|c}
                           &   Property Inference~\lgcite{26} &   DLG   \\ \hline
              Eyeglasses   &    0.94              &   \textbf{1.00}   \\ \hline
              Asian        &    0.92              &   \textbf{1.00}   \\
        \end{tabular}
    \captionof{table}{AUC score on LFW dataset.}
    \label{tab:auc_score}
    }
\end{minipage}
\vspace{-10pt}

\revquote{R3}{Comparison with previous work.} To the best of our acknowledge, \textbf{DLG is the first algorithm that performs pixel-level and token-level leakage} based on shared gradients. We have compared conventional synthetic outputs and our recovered results in the Fig.~4 in our paper. We also add a comparison on property inference task in Tab.~\ref{tab:auc_score}: DLG is significantly better since it can directly obtain the raw training data. In the revision, we will add more comparisons.


% common questions
\revquote{R1, R2}{Trade-off between accuracy and defendability.} 
\revquote{R3}{The method is easy to defend.} 
\revquote{R1}{Does 8-bit help?}
DLG is not easy to defend unless with a significant drop in accuracy. 8-bit gradient does not help either. 
We study the trade-off between accuracy and defendability in Tab.~\ref{tab:defense_accuracy}. It shows that \textbf{only when the defense strategy starts to degrade the accuracy then the deep leakage can be defended}. 
% G: Gaussian noise, L: Laplacian noise and also lower precision. 

\begin{table}[h!]
\setlength{\tabcolsep}{4.5pt}
\centering
\footnotesize{
\begin{tabular}[t]{c||c||c|c|c|c||c|c|c|c||c|c}
    & Original  & \textbf{G}-$10^{-4}$ & \textbf{G}-$10^{-3}$ & \textbf{G}-$10^{-2}$ & \textbf{G}-$10^{-1}$ & \textbf{L}-$10^{-4}$ & \textbf{L}-$10^{-3}$ & \textbf{L}-$10^{-2}$ & \textbf{L}-$10^{-1}$  & \textbf{FP}-16 & \textbf{8 bit}  \\ \hline
    Accuracy & 76.3\% & 75.6\% & 73.3\% & 45.3\% & $\le$1\% &  75.6\%  &  73.4\%    &  46.2\%   &  $\le$1\%  & 76.1\%  & 53.7\% \\ \hline
    Defendability & -- & \lgxmark &  \lgxmark & \lgcmark &  \lgcmark & \lgxmark & \lgxmark & \lgcmark & \lgcmark & \lgxmark & \lgcmark
\end{tabular}
}
\caption{\textbf{G}: Gaussian noise, \textbf{L}: Laplacian noise, \textbf{FP}: Floating number. 
\lgcmark~means it successfully defends against DLG while \lgxmark~means fails to defend. The accuracy is evaluated on CIFAR-100, same as what we used in the paper.
}
\label{tab:defense_accuracy}
\vspace{-15pt}
\end{table}

\revquote{R2}{Do you use all trainable parameters of the ResNet as $\nabla W$? Is the model $W$ trained to convergence?} 
Gradients of \textit{all} trainable parameters are used as $\nabla W$. It is important to clarify that DLG does \textit{not} require the model trained to converge:~\textbf{The attack can be performed at any moment during the training}~(Tab.~\ref{tab:leak_vs_convergence}). Our results in paper are based on randomly initialized models. 

\begin{table}[h!]
    \centering
    \small{
    \begin{tabular}{c|c|c|c|c}
        Train Progress & 0\% epochs & 30\% epochs & 70\% epochs & 100\% epochs \\ \hline
         MSE      & $5.7 \times 10^{-6}$ & $3.1 \times 10^{-7}$ & $4.4 \times 10^{-6}$ &  $3.3 \times 10^{-6}$
    \end{tabular}
    \caption{The MSE between leaked image and ground truth on different training stages. Pixel values are normalized to $[0, 1]$. The leaked image is nearly identical to original ones at each training phase.}
    \label{tab:leak_vs_convergence}
    }
\end{table}
\vspace{-15pt}

\revquote{R1}{The concept of `iterations' } % \hspace{-5pt}
refers to the ``n'' in the for-loop in DLG algorithm, not the training iterations.

% \revquote{R1}{8 bits precision, for example} \revquote{R1}{F is twice differentiable, any other conditions?}

\revquote{R2}{Fig 5. Is the blue line "L2 distance" over all parameters and other lines over parameters of specific layers?}
No, the L2 distance (on the top of the figure) is measured between the leaked image and original ground truth image. Other lines are the distance between dummy gradients~$\nabla W'$ and real gradients~$\nabla W$ in each layer. We'll make it clear in the paper.

\revquote{R2}{Are qualitative results from a held-out test set? How/why were these images chosen?} Yes, qualitative results are from a held-out test set. These images are randomly sampled and more examples have been already provided in the appendix. There is no cherry-picking. 

\revquote{R2}{DLG becomes harder (needs more iterations) to attack when batch size increases." Wouldn't this also make for a good defense? How do the reconstruction results vary with batch size?}
In multi-batch examples (Fig 3 in paper and Line 1 in appendix), we only observe few artifact pixels compared with single-batch cases. Note DLG can be performed off-line as long as current model status and gradients are saved. Large-batch is not a good defense strategy since the information can be still leaked with more time and iterations.


\revquote{R3}{Time cost for L-BFGS reconstruction.} 
Though L-BFGS takes more calculations for every single step, it is still faster (5~minutes) than other optimizers like SGD (30~minutes) on our hardware (Nvidia Tesla v100).


\begingroup
\footnotesize{
\renewcommand{\section}[2]{} % omit the section title 
\bibliographystyle{abbrv}
\bibliography{references}
}
\endgroup

\end{document}